\begin{abstract}
We present \textbf{WiFi-TTA-Bench}, a harm-aware benchmark for test-time adaptation (TTA) on WiFi Channel State Information (CSI) under \emph{physics-structured} domain shift. Our headline finding is a clean counterexample against the ``TTA always helps'' and the ``TTA always fails'' narratives: the \emph{same} BatchNorm-equipped MLP+BN backbone with faithful BN-only TENT achieves $+$21.1~pp on SignFi's milder temporal split but $-$0.61~pp on Widar\_BVP's harder cross-room shift ($n{=}15$ paired observations each). This suggests that on the evaluated WiFi CSI splits, TTA outcomes depend on the interaction between shift severity and architecture rather than on the choice of adaptation method alone. On the primary MLP backbone, no evaluated method's 95\% bootstrap CI clears source-only inference on any of the three real datasets, and our simplified shared-MLP ports of SAR and CoTTA are strongly negative across all three datasets (Widar $-$7.6/$-$6.6, NTU-Fi $-$17.7/$-$17.7, SignFi $-$11.8/$-$12.4~pp), while faithful BN-only SAR/CoTTA ports on an MLP+BN Widar backbone shrink to $-$0.15/$-$0.38~pp with CIs overlapping zero. The benchmark covers up to 12 methods on Widar\_BVP (10--12 on the other datasets; coverage matrix in Appendix~\ref{app:coverage}), with explicit fidelity labels (TENT-proj vs faithful BN-only TENT; SHOT as IM; simplified SAR/CoTTA/LAME ports). We present evidence consistent with three failure mechanisms (overconfident target predictions with mean max-softmax confidence 70--81\,\% across every class despite per-class accuracies as low as 24\,\%; batch-statistic unreliability at WiFi-scale batches; and a large synthetic-to-real swing for the physics-prior objective family), and show that source calibration via label smoothing reduces \texttt{physics\_tta} damage on Widar (not significant at $n{=}15$). WiFi-TTA-Bench delivers standardized splits, machine-readable metadata (Croissant core + RAI fields), a harm-aware 5-metric evaluation protocol (gain, CI, Cohen's $d$, negative-adaptation rate, source drop), and a tested artifact (248 automated tests) with component scripts that regenerate every main table and figure from prepared data.
\end{abstract}
